\section{评估策略与实验设计}

\subsection{离线指标与对比}

离线评估通常在与训练集不同的 holdout 环境中完成，以确保策略并未过拟合。常见指标包括累积 reward、episode 长度、稳定性（标准差）和 sample efficiency。下面的表格给出了这几个指标的关注点：

\begin{table}[H]
\centering
\begin{tabular}{lp{8cm}}
\toprule
指标 & 备注 \\
\midrule
累计 reward & 反映策略在固定起点下的任务完成度 \\
Episode length & 评估是否学会了更简洁的路径 / 效率 \\
Sample efficiency & 每单位环境交互带来的 reward 改善幅度 \\
稳定性 & 不同 seeds 下 reward 的波动，过大会导致部署风险 \\
\bottomrule
\end{tabular}
\caption{离线评估常用指标}
\end{table}

\begin{importantbox}{Sample efficiency 与 reward ceiling 的权衡}
有时候提升累计 reward 的边际收益非常低，但 sample efficiency 仍然未稳定，说明 agent 还在反复探索。建议在迭代曲线上同时跟踪两者，并在 plateau 期引入更强的 exploration 或 auxiliary task。
\end{importantbox}

\subsection{在线对照与干预}

在生产环境下，可采用 champion/challenger 或 A/B 形式对照测试新策略。应保证 challenger 只在较小流量下运行，并准备回滚策略。

\begin{knowledgebox}{A/B vs Champion/ Challenger}
A/B：线上流量按照比例直接路由到 challenger，适合 behavior policy 与 baseline 差异较小时；Champion/Challenger：在 shadow 环境运行 challenger，通过 offline evaluation 决策是否 promote。
\end{knowledgebox}

\begin{warningbox}{评价分布偏移}
训练环境与部署环境一旦存在观测噪声差异，离线指标的信号就可能失真。在安全敏感系统（如自动驾驶）中要优先建立 scenario 库用于 stress test。
\end{warningbox}

\subsection{本章小结}

本章梳理了 Q-Learning 的评估体系，强调同时关注累积 reward、样本效率与稳定性，并建议在部署前设置对照测试与 stress test。

